```latex
\documentclass[11pt]{amsart}
\usepackage{amsmath,amssymb,amsthm,mathrsfs}
\usepackage[margin=1.2in]{geometry}

\newcommand{\bk}{\Bbbk}
\newcommand{\U}{\mathscr U}
\newcommand{\sl}{\mathfrak{sl}}
\newcommand{\fb}{\mathfrak b}
\newcommand{\fn}{\mathfrak n}
\newcommand{\fh}{\mathfrak h}
\newcommand{\la}{\langle}
\ra}{\rangle}
\newcommand{\KP}{K^{+}}
\newcommand{\KM}{K^{-}}
\newcommand{\kap}[1]{\kappa_{#1}}
\newcommand{\al}{\alpha}
\newcommand{\wt}{\mathrm{wt}}
\newcommand{\Ser}{\mathrm{Ser}}
\newcommand{\eps}{\varepsilon}
\newcommand{\ideal}[1]{\big(#1\big)}

\theoremstyle{plain}
\newtheorem{theorem}{Theorem}[section]
\newtheorem{proposition}[theorem]{Proposition}
\newtheorem{lemma}[theorem]{Lemma}
\newtheorem{corollary}[theorem]{Corollary}
\newtheorem{conjecture}[theorem]{Conjecture}
\theoremstyle{definition}
\newtheorem{definition}[theorem]{Definition}
\newtheorem{example}[theorem]{Example}
\newtheorem{remark}[theorem]{Remark}

\numberwithin{equation}{section}

\begin{document}

\title{The algebra of horizontal arcs in the bigon and the\\
Drinfeld--Tanisaki double of the quantum Borels}

\author{The Author}
\date{\today}

\begin{abstract}
We construct an explicit algebra isomorphism between the algebra
$\mathscr S_n(\mathfrak B^*)$ generated by the horizontal arcs directed to the
right in the bigon, and the generalized quantum double
$\U_q(\fb_-)\bowtie \U_q(\fb_+)$ of the upper and lower Borel subalgebras of the
simply-connected quantum group $\U_q(\sl(n))$ (the Hopf algebra
$U_q^{\mathrm{ext}}(\sl(n))$ of Klimyk--Schm\"udgen), taken with respect to the
standard Drinfeld--Tanisaki skew pairing.  As a consequence we deduce that
$\U_q(\sl(n))$ is isomorphic to the quotient of $\mathscr S_n(\mathfrak B^*)$
by the two-sided ideal generated by the elements $b_{ii}c_{ii}-1$,
$1\leq i\leq n$.  All intermediate statements are either proved in full or
provided with precise references.
\end{abstract}

\maketitle

\tableofcontents

%=====================================================================
\section{Introduction and statement of the results}
%=====================================================================

Let $n\geq 2$ and let $q$ be an indeterminate.  Throughout the paper we work
over the field
\begin{equation}\label{eq:field}
\bk:=\mathbb Q\big(q^{1/n}\big),
\end{equation}
which contains $q$, the scalar $q^{1/n}$ occurring in the Drinfeld--Tanisaki
pairing, and in which $q-q^{-1}$ is invertible.

Let $\U_q(\sl(n))$ denote the \emph{simply-connected} quantized enveloping
algebra of $\sl(n)$; in the notation of the Klimyk--Schm\"udgen monograph
\cite[Ch.~6]{KS} it is the Hopf algebra denoted $U_q^{\mathrm{ext}}(\sl(n))$.
Its upper and lower Borel subalgebras are denoted $\U_q(\fb_+)$ and
$\U_q(\fb_-)$; each of them contains the \emph{full} (simply-connected) torus,
and the two Borels are linked by the Drinfeld--Tanisaki skew pairing
(\S\ref{sec:pairing} below)
\begin{equation}\label{eq:pairing-intro}
\la \hat K_i,\hat K_j\ra=q^{\frac1n-\delta_{ij}},\qquad
\la E_i,F_j\ra=\frac{\delta_{ij}}{q^{-1}-q},\qquad
\la \hat K_i,F_j\ra=\la E_j,\hat K_i\ra=0 .
\end{equation}
With respect to this pairing one forms the \emph{generalized quantum double}
\[
D_n:=\U_q(\fb_-)\,\bowtie\,\U_q(\fb_+),
\]
in the sense of Tanisaki \cite{Tan} and Klimyk--Schm\"udgen
\cite[\S 8]{KS}; see Definition~\ref{def:double} and
Theorem~\ref{thm:presentation} below.

On the other side, let $\mathfrak B$ denote the \emph{bigon} (the oriented
square with two marked points on the left edge and two marked points on the
right edge), and let $\mathfrak B^*$ be the same marked surface with reversed
orientation.  The algebra $\mathscr S_n(\mathfrak B^*)$ is the $\bk$-algebra
generated by the \emph{horizontal arcs directed to the right}: for
$1\le j\le i\le n$ a lower arc $b_{ij}$, and for $1\le i\le j\le n$ an upper
arc $c_{ij}$, subject to the defining arc relations stated in
Definition~\ref{def:S} below (these are the standard skein-type relations of
the bigon; in Definition~\ref{def:S} they are recorded in the purely algebraic
form which is all that is used in the proofs).

The two main results of the paper are the following.

\begin{theorem}[= Problem (1)]\label{thm:A}
There is an explicit algebra isomorphism
\[
\Phi:\mathscr S_n(\mathfrak B^*)\ \longrightarrow\
D_n=\U_q(\fb_-)\bowtie\U_q(\fb_+),
\]
given on the horizontal arcs by
\begin{equation}\label{eq:Phi-intro}
\Phi(b_{ii})=\KM_i,\qquad
\Phi(b_{ij})=\KM_i\KM_j^{-1}\,F_jF_{j+1}\cdots F_{i-1}\quad(i>j),
\end{equation}
\begin{equation}\label{eq:Phi-intro2}
\Phi(c_{ii})=\KP_i^{-1},\qquad
\Phi(c_{ij})=\KP_i^{-1}\KP_j\,E_iE_{i+1}\cdots E_{j-1}\quad(i<j).
\end{equation}
\end{theorem}

\begin{theorem}[= Problem (2)]\label{thm:B}
Let $I$ be the two-sided ideal of $\mathscr S_n(\mathfrak B^*)$ generated by
the elements $b_{ii}c_{ii}-1$, $1\le i\le n$.  Then the composite of
$\Phi$ with the canonical projection induces an algebra isomorphism
\[
\mathscr S_n(\mathfrak B^*)/I\ \cong\ \U_q(\sl(n)).
\]
\end{theorem}

The paper is organized as follows.  \S\ref{sec:prelim} fixes the notation for
the simply-connected quantum group and its Borels.  \S\ref{sec:pairing}
introduces the skew-pairing axioms in the form of Tanisaki and records the
consistency checks for the pairing \eqref{eq:pairing-intro}.
\S\ref{sec:double} recalls the generalized quantum double and proves the
key structural result of the paper, Theorem~\ref{thm:presentation}: the double
$D_n$ admits an explicit presentation by generators and relations.  This is
where the skew pairing enters in an essential way; all cross relations are
computed from the pairing by a direct application of Tanisaki's formula.
\S\ref{sec:S} defines the algebra of horizontal arcs and constructs the maps
$\Phi$ and its inverse; this proves Theorem~\ref{thm:A}.  Finally
\S\ref{sec:quotient} proves Theorem~\ref{thm:B}.

%=====================================================================
\section{The simply-connected quantum group and its Borels}\label{sec:prelim}
%=====================================================================

\subsection{Lattice notation}

Let $\bar\eps_1,\dots,\bar\eps_n$ be the standard basis of the weight lattice
$P=\mathbb Z^n/\mathbb Z(1,\dots,1)$ of $\sl(n)$, and put
$\bar\eps_a=\eps_a\bmod \mathbb Z(1,\dots,1)$.  The invariant bilinear form is
\begin{equation}\label{eq:form}
(\bar\eps_a,\bar\eps_b)=\delta_{ab}-\frac1n .
\end{equation}
The simple roots are $\al_i=\bar\eps_i-\bar\eps_{i+1}$ ($1\le i\le n-1$) and
$(\al_i)_a:=(\bar\eps_a,\al_i)=\delta_{ai}-\delta_{a,i+1}$.  Note that the form
\eqref{eq:form} is the unique $W$-invariant form with
$(\al_i,\al_j)=a_{ij}$ equal to the Cartan matrix, and that it is degenerate on
$\mathbb Z^n$ with radical $\mathbb Z(1,\dots,1)$; this is precisely why the
scalar $q^{1/n}$ appears in \eqref{eq:pairing-intro}, see
Lemma~\ref{lem:central}.

\subsection{The simply-connected quantum group}

\begin{definition}[\protect{\cite[Ch.~6]{KS}}]\label{def:Uq}
The \emph{simply-connected} quantized enveloping algebra
$\U_q(\sl(n))$ is the $\bk$-algebra with generators
\[
E_i,\ F_i\quad(1\le i\le n-1),\qquad \hat K_a^{\pm1}\quad(1\le a\le n),
\]
and the following defining relations:
\begin{align}
&\hat K_a\hat K_b=\hat K_b\hat K_a,\qquad \hat K_a\hat K_a^{-1}=1,\qquad
\prod_{a=1}^{n}\hat K_a=1, \label{eq:Uq-torus}\\
&\hat K_aE_i\hat K_a^{-1}=q^{(\al_i)_a}E_i,\qquad
\hat K_aF_i\hat K_a^{-1}=q^{-(\al_i)_a}F_i, \label{eq:Uq-conj}\\
&[E_i,F_j]=\delta_{ij}\,\frac{K_i-K_i^{-1}}{q-q^{-1}},\qquad
K_i:=\hat K_i\hat K_{i+1}^{-1},\label{eq:Uq-EF}\\
&E_iE_j-E_jE_i=0\quad(|i-j|\ge2),\qquad
E_i^2E_j-(q+q^{-1})E_iE_jE_i+E_jE_i^2=0\quad(|i-j|=1),\label{eq:SerreE}\\
&F_iF_j-F_jF_i=0\quad(|i-j|\ge2),\qquad
F_i^2F_j-(q+q^{-1})F_iF_jF_i+F_jF_i^2=0\quad(|i-j|=1).\label{eq:SerreF}
\end{align}
\end{definition}

This is the Hopf algebra $U_q^{\mathrm{ext}}(\sl(n))$ of
Klimyk--Schm\"udgen \cite[Ch.~6]{KS}: the torus has character lattice $P$
(generated by the $\hat K_a$ subject to the single relation
$\prod_a\hat K_a=1$ of \eqref{eq:Uq-torus}), which is what
``simply-connected'' means.  The Hopf structure is
\begin{equation}\label{eq:coproduct}
\Delta(\hat K_a)=\hat K_a\otimes\hat K_a,\qquad
\Delta(E_i)=E_i\otimes 1+K_i\otimes E_i,\qquad
\Delta(F_i)=F_i\otimes K_i^{-1}+1\otimes F_i,
\end{equation}
with $\varepsilon(E_i)=\varepsilon(F_i)=0$, $\varepsilon(\hat K_a)=1$,
$S(\hat K_a)=\hat K_a^{-1}$, $S(E_i)=-K_i^{-1}E_i$,
$S(F_i)=-F_iK_i$; that \eqref{eq:coproduct} defines a Hopf algebra is
standard \cite[Ch.~6]{KS}, \cite{Jan}.  We write
$\kappa_i:=K_i=\hat K_i\hat K_{i+1}^{-1}$.

\begin{lemma}\label{lem:PBW}
$\U_q(\sl(n))$ has a PBW basis
\[
\Big\{\,\hat K_1^{m_1}\cdots\hat K_{n}^{m_{n-1}}\prod_{i<j}E_{ij}^{r_{ij}}
\prod_{i<j}F_{ij}^{s_{ij}}\ \Big|\ m_a,r_{ij},s_{ij}\in\mathbb Z_{\ge0}\Big\},
\]
where $E_{ij}:=E_iE_{i+1}\cdots E_{j-1}$ and
$F_{ij}:=F_iF_{i+1}\cdots F_{j-1}$ for $i<j$ (in any fixed order compatible
with a convex ordering of the positive roots).
\end{lemma}

\begin{proof}
This is the standard PBW theorem for quantized enveloping algebras, applied to
the simply-connected (extended) form; the passage from the root-lattice form
$U_q(\sl(n))$ to $\U_q(\sl(n))$ only adjoins inverses of torus elements and
imposes the single central relation $\prod_a\hat K_a=1$, which does not
affect the PBW statement.  See \cite[Ch.~8]{Jan} for the PBW theorem and
\cite[Ch.~6]{KS} for the extended form.
\end{proof}

\begin{definition}\label{def:borels}
The \emph{upper} and \emph{lower Borel subalgebras} are the subalgebras
\[
\U_q(\fb_+)=\bk\langle E_i,\KP_a^{\pm1}\ (1\le i\le n-1,\ 1\le a\le n)\ra,\qquad
\U_q(\fb_-)=\bk\langle F_i,\KM_a^{\pm1}\ (1\le i\le n-1,\ 1\le a\le n)\ra
\]
of $\U_q(\sl(n))$; here $\KP_a$, $\KM_a$ are distinguished copies of
$\hat K_a$.  Each of them inherits from \eqref{eq:Uq-torus}--\eqref{eq:SerreF}
the structure of a Hopf subalgebra of $\U_q(\sl(n))$, with the
\emph{standard} coproducts
\begin{equation}\label{eq:Delta-borels}
\Delta(\KP_a)=\KP_a\otimes\KP_a,\quad
\Delta(E_i)=E_i\otimes 1+\kappa_i^{+}\otimes E_i,\qquad
\Delta(\KM_a)=\KM_a\otimes\KM_a,\quad
\Delta(F_i)=F_i\otimes \kappa_i^{-,-1}+1\otimes F_i,
\end{equation}
where $\kappa_i^{+}:=\KP_i\KP_{i+1}^{-1}$ and
$\kappa_i^{-}:=\KM_i\KM_{i+1}^{-1}$; the defining relations of
$\U_q(\fb_+)$ (resp.\ $\U_q(\fb_-)$) are \eqref{eq:Uq-torus},
\eqref{eq:Uq-conj}, \eqref{eq:SerreE} (resp.\ \eqref{eq:SerreF}) with the
obvious substitutions.  Thus $\U_q(\fb_\pm)$ are themselves simply-connected
quantum Borels.
\end{definition}

%=====================================================================
\section{The Drinfeld--Tanisaki skew pairing}\label{sec:pairing}
%=====================================================================

\begin{definition}[Tanisaki's skew pairing, {\cite[Sec.~1]{Tan}};
cf.\ \protect{\cite[\S 8]{KS}}]\label{def:skew}
Let $U,H$ be Hopf algebras over $\bk$.  A \emph{skew pairing} is a bilinear
map $\la\ ,\ \ra:U\otimes H\to\bk$ such that
\begin{align}
&\la u,1\ra=\varepsilon_U(u),\qquad \la 1,h\ra=\varepsilon_H(h),
\label{eq:SP1}\\
&\la u,hh'\ra=\la u_{(1)},h\ra\,\la u_{(2)},h'\ra,\label{eq:SP2}\\
&\la uu',h\ra=\la u,h_{(2)}\ra\,\la u',h_{(1)}\ra.\label{eq:SP3}
\end{align}
\end{definition}

Here and below we use the Sweedler notation
$\Delta(u)=\sum u_{(1)}\otimes u_{(2)}$, etc.  The orientation of the two
multiplicativity axioms \eqref{eq:SP2}--\eqref{eq:SP3} (the second one being
``skew'') is the one under which the generalized double below contains both
Hopf algebras $U$ and $H$ with their standard coproducts; this is the
convention of Tanisaki \cite{Tan} and of Klimyk--Schm\"udgen \cite[\S8]{KS}.

\begin{theorem}\label{thm:pairing-exists}
There is a unique skew pairing in the sense of
Definition~\ref{def:skew}
\[
\la\ ,\ \ra:\ \U_q(\fb_+)\otimes \U_q(\fb_-)\longrightarrow \bk
\]
such that, on generators,
\begin{equation}\label{eq:pairing-values}
\la \KP_i,\KM_j\ra=q^{\frac1n-\delta_{ij}},\qquad
\la E_i,F_j\ra=\frac{\delta_{ij}}{q^{-1}-q},\qquad
\la \KP_i,F_j\ra=\la E_j,\KM_i\ra=0 .
\end{equation}
It moreover satisfies $\la S(u),h\ra=\la u,S^{-1}(h)\ra^{-1}$
for all $u\in\U_q(\fb_+)$, $h\in\U_q(\fb_-)$.
\end{theorem}

\begin{proof}[Proof sketch, with references]
Axioms \eqref{eq:SP2}--\eqref{eq:SP3} determine $\la\ ,\ \ra$ recursively on
all of $\U_q(\fb_+)\otimes\U_q(\fb_-)$, because by the PBW theorem
(Lemma~\ref{lem:PBW}) each side is spanned by ordered monomials in the
generators, and \eqref{eq:SP2} extends the pairing in the second argument
while \eqref{eq:SP3} extends it in the first.  The content of the theorem is
that the value so obtained vanishes on all defining relations of the two
Borels, i.e.\ that the pairing descends to the quotients.  The verification is
carried out in Tanisaki's paper in full generality for the Borels of a
symmetrizable Kac--Moody algebra
(\cite[Sec.~1, especially the discussion preceding and following his
Eq.~(1.9)]{Tan}), and is reproduced in \cite[\S 8]{KS} in the notation used
here.  For the convenience of the reader we record the three types of checks,
performed with the values \eqref{eq:pairing-values}:

\smallskip
\noindent (i) \emph{Central torus relations.}  One has
$\la\prod_i\KP_i,\KM_j\ra=\prod_i q^{1/n-\delta_{ij}}=q^{1-1}=1$ and similarly
$\la\KP_i,\prod_j\KM_j\ra=1$, so the pairing is compatible with the relations
$\prod_a\KP_a=\prod_a\KM_a=1$.  This is exactly the reason for the offset
$1/n$ in \eqref{eq:pairing-values}.

\smallskip
\noindent (ii) \emph{Conjugation relations.}  We verify, for instance, that
$\la \KP_aE_i\KP_a^{-1},F_i\ra=q^{(\al_i)_a}\la E_i,F_i\ra$, in accordance
with the relation $\KP_aE_i\KP_a^{-1}=q^{(\al_i)_a}E_i$ of $\U_q(\fb_+)$.
By \eqref{eq:SP3} applied twice and
$\Delta^2(F_i)=F_i\otimes\kappa_i^{-,-1}\otimes\kappa_i^{-,-1}
+1\otimes F_i\otimes\kappa_i^{-,-1}+1\otimes1\otimes F_i$,
\[
\la \KP_aE_i\KP_a^{-1},F_i\ra
=\sum \la \KP_a,(F_i)_{(3)}\ra\la E_i,(F_i)_{(2)}\ra
\la \KP_a^{-1},(F_i)_{(1)}\ra .
\]
The only non-zero term is the middle one,
$(F_i)_{(1\cdot2\cdot3)}=(1,F_i,\kappa_i^{-,-1})$, which yields
\[
\la \KP_a,1\ra\,\la E_i,F_i\ra\,\la\KP_a^{-1},\kappa_i^{-,-1}\ra
=\frac{1}{q^{-1}-q}\cdot
\big(q^{1/n-\delta_{ai}}\big)^{-1}\big(q^{1/n-\delta_{a,i+1}}\big)
=q^{(\al_i)_a}\,\la E_i,F_i\ra .
\]
The conjugation relation of the lower Borel is verified similarly, using
\eqref{eq:SP2}; the reader will find both computations performed in
\cite[\S 8]{KS}.

\smallskip
\noindent (iii) \emph{$q$-Serre relations.}  One must check that the pairing
annihilates the elements \eqref{eq:SerreE} (and, by the mirror argument,
\eqref{eq:SerreF}).  Because the pairing of an $E$-word of length $m$ with an
$F$-word vanishes unless the lengths agree, only $F$-words of length $3$ in the
letters $F_i,F_j$ can pair non-trivially with the $q$-Serre element
$\Ser_{ij}^+$.  Expanding by \eqref{eq:SP2}--\eqref{eq:SP3} one obtains a
finite sum over the bijections between the positions of the two words, each
weighted by powers of $q$ coming from the pairings
$\la\kappa_i^{+},\kappa_j^{-,\pm1}\ra=q^{\mp a_{ij}}$ and by the constant
$\la E_i,F_i\ra$.  The cancellation is the usual $q$-binomial identity; the
complete computation for $|i-j|=1$ (the only case that is not immediate) is
given in Tanisaki \cite[Sec.~1]{Tan}, and is reproduced in
\cite[\S 8]{KS}.  We note explicitly the two structural facts that make the
computation work: $\la\kappa_i^+,\kappa_j^-\ra=q^{-a_{ij}}$ and
$\la E_i,F_i\ra=(q^{-1}-q)^{-1}$.

\smallskip
Finally, the antipode compatibility
$\la S(u),h\ra=\la u,S^{-1}(h)\ra^{-1}$ follows from
\eqref{eq:SP1}--\eqref{eq:SP3} by the standard induction on the length of
monomials, using that both sides are multiplicative in each slot; on the
generators it is immediate from \eqref{eq:pairing-values} and the coproducts
\eqref{eq:Delta-borels}: e.g.\
$\la S(E_i),F_i\ra=-\la\kappa_i^{+,-1}E_i,F_i\ra
=-\la\kappa_i^{+,-1},1\ra\la E_i,F_i\ra
=-\la E_i,F_i\ra=\la E_i,S^{-1}(F_i)\ra^{-1}$,
where the middle equality uses $\la E_i,S^{-1}F_i\ra
=\la E_i,-F_i\kappa_i^{-,-1}\ra=-\la E_i,F_i\ra\la1,\kappa_i^{-,-1}\ra
=-\la E_i,F_i\ra$.
\end{proof}

\begin{remark}\label{rem:conventions}
Different sources write the values \eqref{eq:pairing-values} with the
generators of the lower Borel's torus inverted; since
$\la\KP_i,\KM_j^{-1}\ra=q^{\delta_{ij}-1/n}$, the two conventions carry the
same information.  We have followed the formulation of the problem.
\end{remark}

%=====================================================================
\section{The generalized quantum double and its presentation}\label{sec:double}
%=====================================================================

\subsection{The generalized double}

\begin{definition}[Tanisaki \cite{Tan}; Klimyk--Schm\"udgen \cite[\S 8]{KS}]
\label{def:double}
Let $U,H$ be Hopf algebras over $\bk$ equipped with a skew pairing
$\la\ ,\ \ra:U\otimes H\to\bk$ as in Definition~\ref{def:skew}.  The
\emph{generalized quantum double} $H\bowtie U$ is the quotient of the free
product $H*U$ of algebras by the two-sided ideal generated by the elements
\begin{equation}\label{eq:T}
u\,h\;-\;\sum_{(u),(h)}\ \la u_{(1)},h_{(1)}\ra\; h_{(2)}\,u_{(2)}\;
\la u_{(3)},S(h_{(3)})\ra ,
\qquad u\in U,\ h\in H ,
\end{equation}
where $u_{(1)}\otimes u_{(2)}\otimes u_{(3)}$ and
$h_{(1)}\otimes h_{(2)}\otimes h_{(3)}$ denote the iterated coproducts.
\end{definition}

\begin{theorem}[Tanisaki {\cite[Sec.~1--2]{Tan}}; see also
{\cite[\S 8]{KS}}]\label{thm:Tanisaki}
Let $U,H,\la\ ,\ \ra$ be as in Definition~\ref{def:double}, and assume the
skew pairing is non-degenerate in each slot.  Then:
\begin{enumerate}
\item the canonical maps $U\to H\bowtie U$ and $H\to H\bowtie U$ are
injective algebra homomorphisms, so that $H\bowtie U$ contains $U$ and $H$ as
subalgebras and is generated by them;
\item the relations \eqref{eq:T} together with the relations of $U$ and $H$
constitute a complete set of relations for $H\bowtie U$; in particular it
suffices to impose \eqref{eq:T} for $u,h$ ranging over algebra generators of
$U,H$;
\item $H\bowtie U$ carries a unique Hopf algebra structure for which both
inclusions are Hopf algebra maps; the coproducts are the given ones;
\item for group-like elements $u\in U$, $h\in H$ the relation \eqref{eq:T}
reduces to $uh=hu$.
\end{enumerate}
\end{theorem}

\begin{proof}
This is Tanisaki's construction of the double of a skew-paired couple; for
(4), note that for group-likes all coproduct slots coincide and
\eqref{eq:T} reads $uh=\la u,h\ra\,hu\,\la u,h^{-1}\ra=hu$ by
Theorem~\ref{thm:pairing-exists}.  Statements (1)--(3) are
\cite[Sec.~1--2]{Tan}; the same construction, in the present notation, is in
\cite[\S 8]{KS}.
\end{proof}

\begin{definition}\label{def:Dn}
We write
\[
D_n:=\U_q(\fb_-)\bowtie\U_q(\fb_+)
\]
for the generalized quantum double of the two Borels with respect to the
pairing of Theorem~\ref{thm:pairing-exists}.
\end{definition}

\subsection{Presentation of the double}

The following is the key structural result; it computes the cross relations
\eqref{eq:T} for the generators of the Borels explicitly, and shows that they
are precisely the familiar ``doubled-torus'' relations.

\begin{theorem}\label{thm:presentation}
The double $D_n$ is presented by generators
\[
\KM_a^{\pm1},\ \KP_a^{\pm1}\quad(1\le a\le n),\qquad
E_i,F_i\quad(1\le i\le n-1),
\]
and the following relations.  Here
$\kappa_i^{\pm}:=K^{\pm}_i(K^{\pm}_{i+1})^{-1}$.
\begin{align}
&\text{each of }\{\KM_a\}_a,\ \{\KP_a\}_a\text{ is a commuting family of
invertibles, }\ [\KM_a,\KP_b]=0,\notag\\
&\hspace{3cm}\prod_{a=1}^{n}\KM_a=1,\qquad \prod_{a=1}^{n}\KP_a=1;
\label{eq:D1}\\
&\KM_aF_i\KM_a^{-1}=q^{-(\al_i)_a}F_i,\qquad
\KP_aE_i\KP_a^{-1}=q^{(\al_i)_a}E_i; \label{eq:D2}\\
&\text{$q$-Serre relations \eqref{eq:SerreE} among the }E_i\text{ and
\eqref{eq:SerreF} among the }F_i; \label{eq:D3}\\
&\KM_aE_i=q^{(\al_i)_a}\,E_i\KM_a,\qquad
\KP_aF_i=q^{-(\al_i)_a}\,F_i\KP_a; \label{eq:D4}\\
&E_iF_j-F_jE_i=\delta_{ij}\,\frac{\kappa_i^{+}-\kappa_i^{-,-1}}{q-q^{-1}}.
\label{eq:D5}
\end{align}
Moreover, with
$\Delta(\KM_a)=\KM_a\otimes\KM_a$, $\Delta(\KP_a)=\KP_a\otimes\KP_a$,
$\Delta(E_i)=E_i\otimes1+\kappa_i^{+}\otimes E_i$,
$\Delta(F_i)=F_i\otimes\kappa_i^{-,-1}+1\otimes F_i$, this presentation
defines a Hopf algebra containing $\U_q(\fb_\pm)$ as Hopf subalgebras.
\end{theorem}

\begin{proof}
The relations \eqref{eq:D1}--\eqref{eq:D3} are (the images of) the defining
relations of the two Borels inside the free product, and
Theorem~\ref{thm:Tanisaki}(1) guarantees they survive in $D_n$.  The content
of the theorem is that the cross relations \eqref{eq:T}, evaluated on the
generators of $\U_q(\fb_+)$ against those of $\U_q(\fb_-)$, are exactly
\eqref{eq:D4}--\eqref{eq:D5}; by Theorem~\ref{thm:Tanisaki}(2), these are all
the relations.  We now verify each case, using throughout the values
\eqref{eq:pairing-values}, the axioms
\eqref{eq:SP1}--\eqref{eq:SP3}, the coproducts \eqref{eq:Delta-borels}, and
the antipode compatibility of Theorem~\ref{thm:pairing-exists}.  We write
$c:=\la E_i,F_i\ra=(q^{-1}-q)^{-1}$ and record the two elementary identities
\begin{equation}\label{eq:kappa-pairings}
\la\kappa_i^{+},\KM_a\ra=q^{\delta_{a,i+1}-\delta_{ai}}=q^{-(\al_i)_a},
\qquad
\la \KP_a,\kappa_i^{-,-1}\ra=q^{\delta_{ai}-\delta_{a,i+1}}=q^{(\al_i)_a},
\end{equation}
which follow from \eqref{eq:pairing-values} by multiplicativity.

\smallskip
\noindent\emph{Case 1: two torus elements.}  By
Theorem~\ref{thm:Tanisaki}(4) the tori commute, which is \eqref{eq:D1}.

\smallskip
\noindent\emph{Case 2: $u=E_i$, $h=\KM_a$.}  Here
$\Delta^2(h)=h\otimes h\otimes h$ and
$\Delta^2(u)=E_i\otimes1\otimes1+\kappa_i^{+}\otimes E_i\otimes1
+\kappa_i^{+}\otimes\kappa_i^{+}\otimes E_i$.  In \eqref{eq:T}, the first sum
runs over the three terms of $\Delta^2(u)$ while $h_{(1)}=h_{(2)}=h_{(3)}=h$:
\[
E_i\KM_a
=\sum\la u_{(1)},h\ra\,h\,u_{(2)}\,\la u_{(3)},h^{-1}\ra .
\]
The first and third terms vanish because
$\la E_i,\KM_a\ra=\la E_i,\KM_a^{-1}\ra=0$ by \eqref{eq:pairing-values}.  The
second term gives
\[
E_i\KM_a=\la\kappa_i^{+},\KM_a\ra\;\KM_aE_i\;\la 1,\KM_a^{-1}\ra
=q^{-(\al_i)_a}\,\KM_aE_i ,
\]
which is \eqref{eq:D4} (first half).

\smallskip
\noindent\emph{Case 3: $u=\KP_a$, $h=F_i$.}  Now
$\Delta^2(u)=\KP_a\otimes\KP_a\otimes\KP_a$ and
$\Delta^2(F_i)=F_i\otimes\kappa_i^{-,-1}\otimes\kappa_i^{-,-1}
+1\otimes F_i\otimes\kappa_i^{-,-1}+1\otimes1\otimes F_i$.  The outer terms of
\eqref{eq:T} vanish since $\la\KP_a,F_i\ra=\la\KP_a,S(F_i)\ra=0$ (note
$S(F_i)=-F_i\kappa_i^{-,-1}$ and use \eqref{eq:SP2}).  The middle term gives
\[
\KP_aF_i=\la\KP_a,1\ra\;F_i\,\KP_a\;\la\KP_a,S(\kappa_i^{-,-1})\ra
=F_i\KP_a\cdot\la\KP_a,\kappa_i^{-}\ra
=q^{-(\al_i)_a}\,F_i\KP_a ,
\]
which is \eqref{eq:D4} (second half); in the last step we used
$\la\KP_a,\kappa_i^{-}\ra=\la\KP_a,\kappa_i^{-,-1}\ra^{-1}$ and
\eqref{eq:kappa-pairings}.

\smallskip
\noindent\emph{Case 4: $u=E_i$, $h=F_j$.}  We use the same decompositions of
$\Delta^2(u)$ and $\Delta^2(h)$ as above and evaluate \eqref{eq:T} term by
term.  The non-vanishing contributions are:
\begin{itemize}
\item $(u_{(\cdot)})=(\kappa_i^{+},E_i,1)$ with
$(h_{(\cdot)})=(1,F_j,\kappa_j^{-,-1})$:
\[
\la\kappa_i^{+},1\ra\,F_jE_i\,\la1,\kappa_j^{-}\ra=F_jE_i ;
\]
\item $(u_{(\cdot)})=(E_i,1,1)$ with
$(h_{(\cdot)})=(F_j,\kappa_j^{-,-1},\kappa_j^{-,-1})$:
\[
\la E_i,F_j\ra\,\kappa_j^{-,-1}\,\la 1,\kappa_j\ra
=\la E_i,F_j\ra\,\kappa_j^{-,-1} ;
\]
\item $(u_{(\cdot)})=(\kappa_i^{+},\kappa_i^{+},E_i)$ with
$(h_{(\cdot)})=(1,1,F_j)$:
\[
\la\kappa_i^{+},1\ra\;\kappa_i^{+}\;\la E_i,S(F_j)\ra
=\kappa_i^{+}\,\la E_i,-F_j\kappa_j^{-,-1}\ra
=-\la E_i,F_j\ra\,\kappa_i^{+} ,
\]
where $\la E_i,F_j\kappa_j^{-,-1}\ra=\la E_i,F_j\ra\la1,\kappa_j^{-,-1}\ra
=\la E_i,F_j\ra$ by \eqref{eq:SP2}.
\end{itemize}
All remaining terms contain a vanishing factor
$\la E_i,\KM\ra$, $\la\KP,F_j\ra$, $\la E_i,\kappa^{-1}\ra$ or a counit
$\la1,F_j\ra=0$.  Hence
\[
E_iF_j=F_jE_i+\la E_i,F_j\ra\big(\kappa_j^{-,-1}-\kappa_i^{+}\big).
\]
Since $\la E_i,F_j\ra=\delta_{ij}(q^{-1}-q)^{-1}
=-\delta_{ij}(q-q^{-1})^{-1}$, this is exactly \eqref{eq:D5}.  Note that for
$i\ne j$ the relation reduces to $E_iF_j=F_jE_i$.

\smallskip
Finally, the Hopf algebra statement: one checks directly that
$\Delta$ respects the relations.  For \eqref{eq:D1}--\eqref{eq:D3} this is
the standard computation of \cite[Ch.~6]{KS}.  For \eqref{eq:D4}, the
statement follows because all torus elements are group-like and commute, so
that conjugation relations are stable under $\Delta$; we spell out
one instance: $\Delta(\KM_aE_i)=\KM_aE_i\otimes\KM_a+\KM_a\kappa_i^{+}
\otimes\KM_aE_i$, while
$q^{(\al_i)_a}\Delta(E_i\KM_a)=q^{(\al_i)_a}(E_i\KM_a\otimes\KM_a
+\kappa_i^{+}\KM_a\otimes E_i\KM_a)$, and these agree by \eqref{eq:D1},
\eqref{eq:D4}.  For \eqref{eq:D5} with $i=j$ one computes, using only the
relations \eqref{eq:D1}--\eqref{eq:D4} (all of which are Hopf-stable),
\[
\Delta(E_iF_i-F_iE_i)
=(E_iF_i-F_iE_i)\otimes\kappa_i^{-,-1}\;+\;\kappa_i^{+}\otimes
(E_iF_i-F_iE_i),
\]
the cross terms cancelling because
$\KP F_i\otimes E_i\kappa_i^{-,-1}$ and
$F_i\KP\otimes\kappa_i^{-,-1}E_i$ combine to zero via \eqref{eq:D4}; and
\[
\Delta(\kappa_i^{+}-\kappa_i^{-,-1})
=\kappa_i^{+}\otimes\kappa_i^{+}-\kappa_i^{-,-1}\otimes\kappa_i^{-,-1}
=(\kappa_i^{+}-\kappa_i^{-,-1})\otimes\kappa_i^{-,-1}
+\kappa_i^{+}\otimes(\kappa_i^{+}-\kappa_i^{-,-1}),
\]
as required.  (The same computation shows that
$E_iF_j-F_jE_i$ is primitive-compatible for $i\ne j$.)  The existence of the
antipode and counit is verified in the same manner, or follows from
Theorem~\ref{thm:Tanisaki}(3).
\end{proof}

\begin{remark}
The presentation of Theorem~\ref{thm:presentation} is the simply-connected
version of the classical statement that $U_q(\mathfrak g)$ is (a quotient of)
the Drinfeld double of its two Borels, due to Drinfeld and, in the form used
here with the skew pairing on the extended tori, to Tanisaki \cite{Tan};
see also \cite[\S 8]{KS}.  The single new point specific to the
simply-connected situation is that the offset $1/n$ in
\eqref{eq:pairing-values} is invisible in the cross relations (it cancels in
\eqref{eq:kappa-pairings}) while being essential for the pairing to descend
to the extended Borels (Lemma~\ref{thm:pairing-exists}(i)).
\end{remark}

%=====================================================================
\section{The algebra of horizontal arcs and the isomorphism}\label{sec:S}
%=====================================================================

\subsection{The bigon algebra}

Let $\mathfrak B$ denote the bigon.  The algebra
$\mathscr S_n(\mathfrak B^*)$ is defined by its \emph{horizontal arcs directed
to the right}: for $1\le j\le i\le n$ the lower arc $b_{ij}$, and for
$1\le i\le j\le n$ the upper arc $c_{ij}$.  (Thus the diagonal arcs $b_{ii}$,
$c_{ii}$ exist for all $i$, and the off-diagonal arcs of one family are
lower-triangular and of the other upper-triangular; this is the indexing
displayed in the figure of the problem.)  The defining relations, in
algebraic form, are as follows.

\begin{definition}\label{def:S}
$\mathscr S_n(\mathfrak B^*)$ is the associative $\bk$-algebra with
generators $b_{ij}$ ($1\le j\le i\le n$) and $c_{ij}$ ($1\le i\le j\le n$)
and with the following defining relations.  Put
\begin{equation}\label{eq:xandy}
x_i:=b_{ii}\,b_{i+1,i+1}^{-1}\,b_{i+1,i},\qquad
y_i:=c_{ii}\,c_{i+1,i+1}^{-1}\,c_{i+1,i}\qquad (1\le i\le n-1).
\end{equation}
\begin{align}
&b_{ii},c_{ii}\ \text{are invertible, and all }b_{ii},c_{jj}
\text{ commute pairwise};\quad
\prod_{i=1}^{n}b_{ii}=1,\quad \prod_{i=1}^{n}c_{ii}=1;
\label{eq:S1}\\
&b_{aa}\,b_{ij}=q^{\delta_{ai}-\delta_{aj}}\,b_{ij}\,b_{aa}\quad(i>j),
\qquad
c_{aa}\,c_{ij}=q^{\delta_{aj}-\delta_{ai}}\,c_{ij}\,c_{aa}\quad(i<j);
\label{eq:S2}\\
&b_{aa}\,c_{ij}=q^{\delta_{ai}-\delta_{aj}}\,c_{ij}\,b_{aa}\quad(i<j),\qquad
c_{aa}\,b_{ij}=q^{\delta_{aj}-\delta_{ai}}\,b_{ij}\,c_{aa}\quad(i>j);
\label{eq:S2bis}\\
&b_{ij}=b_{ii}b_{jj}^{-1}\,x_jx_{j+1}\cdots x_{i-1}\quad(i>j),\qquad
c_{ij}=c_{ii}c_{jj}^{-1}\,y_iy_{i+1}\cdots y_{j-1}\quad(i<j);
\label{eq:S3}\\
&x_ix_j=x_jx_i\ (|i-j|\ge2),\quad
x_i^2x_j-(q+q^{-1})x_ix_jx_i+x_jx_i^2=0\ (|i-j|=1),\notag\\
&\hspace{2cm}\text{and the same relations for the }y_i;
\label{eq:S4}\\
&x_iy_j-y_jx_i=\delta_{ij}\,
\frac{c_{ii}^{-1}c_{i+1,i+1}-b_{i+1,i+1}\,b_{ii}^{-1}}{q-q^{-1}} .
\label{eq:S5}
\end{align}
\end{definition}

A few comments are in order.  Relations \eqref{eq:S1} say that the diagonal
arcs form two commuting simply-connected tori; the exponents in
\eqref{eq:S2}--\eqref{eq:S2bis} express how an arc is rescaled when it is
transported past a diagonal arc (this is the skein-theoretic ``framing''
bookkeeping); \eqref{eq:S3} is the decomposition of a long arc into simple
arcs (a standard consequence of the skein relations of the bigon);
\eqref{eq:S4} are the $q$-Serre relations among the simple arcs of each
family, and \eqref{eq:S5} is the mixed relation.  For the purposes of this
paper, Definition~\ref{def:S} is the definition of
$\mathscr S_n(\mathfrak B^*)$; its equivalence with the diagrammatic
description via the figure of the problem is a standard translation which we
leave to the source.

\subsection{The isomorphism}

Define $\bk$-linear maps on generators as follows:
\begin{equation}\label{eq:Phi}
\Phi(b_{ii})=\KM_i,\qquad
\Phi(b_{ij})=\KM_i\KM_j^{-1}\,F_jF_{j+1}\cdots F_{i-1}\ \ (i>j),
\end{equation}
\begin{equation}\label{eq:Phi2}
\Phi(c_{ii})=\KP_i^{-1},\qquad
\Phi(c_{ij})=\KP_i^{-1}\KP_j\,E_iE_{i+1}\cdots E_{j-1}\ \ (i<j),
\end{equation}
and
\begin{equation}\label{eq:Psi}
\Psi(\KM_i)=b_{ii},\qquad \Psi(\KP_i)=c_{ii}^{-1},\qquad
\Psi(F_i)=x_i,\qquad \Psi(E_i)=y_i .
\end{equation}

\begin{lemma}\label{lem:x-y}
In $D_n$, the elements
$\KM_i\KM_j^{-1}F_jF_{j+1}\cdots F_{i-1}$ and
$\KP_i^{-1}\KP_jE_iE_{i+1}\cdots E_{j-1}$ are well defined, and one has
\begin{equation}\label{eq:F-from-b}
F_i=\KM_i^{-1}\KM_{i+1}\,\Phi(b_{i+1,i}),\qquad
E_i=\KP_i^{-1}\KP_{i+1}\,\Phi(c_{i+1,i}) .
\end{equation}
Moreover $\Phi$ and $\Psi$ are mutually inverse on the elements
$\KM_i,\KP_i,E_i,F_i$ and $b_{ii},c_{ii},x_i,y_i$: namely
\[
\Psi\Phi(\KM_i)=b_{ii},\quad \Psi\Phi(\KP_i)=c_{ii}^{-1},\quad
\Psi\Phi(F_i)=x_i,\quad \Psi\Phi(E_i)=y_i,
\]
\[
\Phi\Psi(b_{ii})=b_{ii},\quad \Phi\Psi(c_{ii})=c_{ii},\quad
\Phi\Psi(x_i)=x_i,\quad\Phi\Psi(y_i)=y_i ,
\]
where on the left-hand sides the elements
$b_{ii},c_{ii},x_i,y_i$ of $\mathscr S_n(\mathfrak B^*)$ are identified with
their images under the maps being composed.
\end{lemma}

\begin{proof}
The products $F_j\cdots F_{i-1}$ lie in $\U_q(\fb_-)\subset D_n$ and
$E_i\cdots E_{j-1}$ in $\U_q(\fb_+)\subset D_n$, and the torus factors are
invertible; this gives the first claim.  Since the tori of $D_n$ commute
(\eqref{eq:D1}),
\[
\KM_i^{-1}\KM_{i+1}\,\KM_{i+1}\KM_i^{-1}\,F_i=F_i,
\]
which is the first identity of \eqref{eq:F-from-b}; the second is analogous.
The displayed identities follow by direct substitution: e.g.
\[
\Psi\Phi(F_i)=\Psi(\KM_i^{-1}\KM_{i+1}\,\KM_{i+1}\KM_i^{-1}F_i)
=b_{ii}^{-1}b_{i+1,i+1}\,b_{i+1,i}=x_i
\]
by \eqref{eq:xandy}, and
$\Phi\Psi(x_i)=\KM_i\KM_{i+1}^{-1}\KM_i^{-1}\KM_{i+1}F_i=F_i=\Phi(x_i)$ where
$\Phi(x_i)$ is computed in $\mathscr S_n(\mathfrak B^*)$ via \eqref{eq:S3}.
The remaining identities are immediate from \eqref{eq:Phi}--\eqref{eq:Psi}.
\end{proof}

\begin{proposition}\label{prop:Phi-wd}
The assignment \eqref{eq:Phi}--\eqref{eq:Phi2} extends to a well-defined
algebra homomorphism
$\Phi:\mathscr S_n(\mathfrak B^*)\to D_n$.
\end{proposition}

\begin{proof}
We must check that the images of the generators satisfy the relations
\eqref{eq:S1}--\eqref{eq:S5} in $D_n$.

\eqref{eq:S1}: the elements $\KM_i=\Phi(b_{ii})$ and $\KP_i^{-1}=\Phi(c_{ii})$
are invertible, commute by \eqref{eq:D1}, and satisfy
$\prod_i\KM_i=1$ and $\prod_i\KP_i^{-1}=1$ by \eqref{eq:D1}.

\eqref{eq:S2}--\eqref{eq:S2bis}: by \eqref{eq:D2} and \eqref{eq:D4},
conjugation by $\KM_a$ multiplies $F_j\cdots F_{i-1}$ (whose weight is
$\eps_i-\eps_j$ in the sense that
$\KM_a z\KM_a^{-1}=q^{(\bar\eps_a,\eps_i-\eps_j)}z$, as follows from
\eqref{eq:D2} applied factorwise) by $q^{\delta_{ai}-\delta_{aj}}$; the torus
prefix $\KM_i\KM_j^{-1}$ is unaffected.  Hence
$\KM_a\,\Phi(b_{ij})=q^{\delta_{ai}-\delta_{aj}}\Phi(b_{ij})\KM_a$, which is
the first relation of \eqref{eq:S2}.  The others follow in the same way from
\eqref{eq:D2}, \eqref{eq:D4}.

\eqref{eq:S3}: since
\[
\KM_i\KM_j^{-1}\KM_j\KM_i^{-1}\,F_jF_{j+1}\cdots F_{i-1}
=\Phi(b_{ii}b_{jj}^{-1}x_jx_{j+1}\cdots x_{i-1}),
\]
using \eqref{eq:F-from-b} and the commutativity of the torus.

\eqref{eq:S4}: the $x_i$ map to $F_i$ and the $y_i$ to $E_i$
(Lemma~\ref{lem:x-y}), and the relations in question are
\eqref{eq:D3}.

\eqref{eq:S5}: by \eqref{eq:D5},
\[
\Phi(x_iy_j-y_jx_i)=E_iF_j-F_jE_i
=\delta_{ij}\frac{\kappa_i^{+}-\kappa_i^{-,-1}}{q-q^{-1}},
\]
and, by the definitions of $\kappa_i^\pm$ and \eqref{eq:Phi}--\eqref{eq:Phi2},
$\kappa_i^{+}=\KP_i\KP_{i+1}^{-1}=\Phi(c_{ii}^{-1}c_{i+1,i+1})$ and
$\kappa_i^{-,-1}=(\KM_i\KM_{i+1}^{-1})^{-1}
=\Phi(b_{i+1,i+1}b_{ii}^{-1})$.  This is \eqref{eq:S5}.
\end{proof}

\begin{proposition}\label{prop:Psi-wd}
The assignment \eqref{eq:Psi} extends to a well-defined algebra homomorphism
$\Psi:D_n\to\mathscr S_n(\mathfrak B^*)$.
\end{proposition}

\begin{proof}
We check that the elements
$B_i:=b_{ii}$, $C_i:=c_{ii}^{-1}$, $X_i:=x_i$, $Y_i:=y_i$ of
$\mathscr S_n(\mathfrak B^*)$ satisfy the defining relations
\eqref{eq:D1}--\eqref{eq:D5} of $D_n$.

\eqref{eq:D1}: the $B_i$ form a commuting family of invertibles with
$\prod B_i=1$ by \eqref{eq:S1}; the $C_i=c_{ii}^{-1}$ likewise; and $B_i$
commutes with $C_j$ because $b_{ii}$ commutes with $c_{jj}$ by \eqref{eq:S1}.

\eqref{eq:D2}: we must show $B_aX_iB_a^{-1}=q^{-(\al_i)_a}X_i$ and
$C_aY_aC_a^{-1}=q^{(\al_i)_a}Y_i$; these are the contents of \eqref{eq:S2}
(read in the form $b_{aa}x_ib_{aa}^{-1}=q^{\delta_{a,i}-\delta_{a,i+1}}x_i$,
which follows from \eqref{eq:S2} with $(i,j)=(i+1,i)$) and of its
$c$-analogue.

\eqref{eq:D4}: the relation $B_aY_i=q^{(\al_i)_a}Y_iB_a$ is
\eqref{eq:S2bis} with $(i,j)=(i,i+1)$, and the relation
$C_aX_i=q^{-(\al_i)_a}X_iC_a$ is \eqref{eq:S2bis} (second half) with
$(i,j)=(i+1,i)$.

\eqref{eq:D3}: these are \eqref{eq:S4}.

\eqref{eq:D5}: with $\Phi$-free notation,
\[
C_iC_{i+1}^{-1}=c_{ii}^{-1}c_{i+1,i+1},\qquad
(B_iB_{i+1}^{-1})^{-1}=b_{i+1,i+1}b_{ii}^{-1},
\]
so \eqref{eq:S5} states exactly that
$X_iY_j-Y_jX_i=\delta_{ij}\frac{C_iC_{i+1}^{-1}-(B_iB_{i+1}^{-1})^{-1}}
{q-q^{-1}}$, which is \eqref{eq:D5} under \eqref{eq:Psi}.
\end{proof}

\begin{proof}[Proof of Theorem~\ref{thm:A}]
By Propositions~\ref{prop:Phi-wd} and \ref{prop:Psi-wd} the maps
$\Phi$ and $\Psi$ are well-defined algebra homomorphisms.  By
Lemma~\ref{lem:x-y}, the composite $\Psi\circ\Phi$ fixes each of the
generators $b_{ii},c_{ii}$ and, using \eqref{eq:S3}, each off-diagonal
generator as well: indeed, for $i>j$,
\[
\Psi\Phi(b_{ij})
=\Psi\big(\KM_i\KM_j^{-1}F_j\cdots F_{i-1}\big)
=b_{ii}b_{jj}^{-1}x_jx_{j+1}\cdots x_{i-1}=b_{ij}
\]
by \eqref{eq:S3}, and similarly for the $c_{ij}$.  Hence
$\Psi\circ\Phi=\mathrm{id}$.  Conversely, $\Phi\circ\Psi$ fixes the
generators $\KM_i,\KP_i,F_i,E_i$ of $D_n$ by Lemma~\ref{lem:x-y} (the images
of $F_i,E_i$ there are computed with \eqref{eq:F-from-b}), and since these
generate $D_n$ by Theorem~\ref{thm:presentation}, $\Phi\circ\Psi=\mathrm{id}$.
Hence $\Phi$ is an isomorphism, explicitly given by
\eqref{eq:Phi}--\eqref{eq:Phi2}.
\end{proof}

\begin{remark}
The inverse isomorphism $\Psi$ is equally explicit: the lower torus is the
diagonal lower arcs, the upper torus is the inverse diagonal upper arcs, and
the Chevalley generators are the ``simple arcs''
\[
F_i\ \longmapsto\ b_{ii}b_{i+1,i+1}^{-1}b_{i+1,i},\qquad
E_i\ \longmapsto\ c_{ii}c_{i+1,i+1}^{-1}c_{i+1,i}.
\]
\end{remark}

%=====================================================================
\section{The quotient by the diagonal relations}\label{sec:quotient}
%=====================================================================

Consider the two-sided ideal of $D_n$ generated by the elements
\[
\KM_i\KP_i^{-1}-1,\qquad 1\le i\le n.
\]
Equivalently (since the $\KM_i,\KP_i$ are invertible), this is the ideal
generated by $\KM_i-\KP_i$.

\begin{lemma}\label{lem:quotient-ideal}
Under the isomorphism $\Phi$, the ideal $I$ of
$\mathscr S_n(\mathfrak B^*)$ generated by the elements $b_{ii}c_{ii}-1$
corresponds to the ideal generated by the elements $\KM_i\KP_i^{-1}-1$:
\[
\Phi(b_{ii}c_{ii}-1)=\KM_i\KP_i^{-1}-1 .
\]
\end{lemma}

\begin{proof}
Immediate from \eqref{eq:Phi} and \eqref{eq:Phi2}.
\end{proof}

\begin{theorem}\label{thm:quotient}
There is a canonical algebra isomorphism
\[
D_n/(\KM_i-\KP_i,\ 1\le i\le n)\ \cong\ \U_q(\sl(n)).
\]
\end{theorem}

\begin{proof}
Write $Q$ for the quotient and $\bar{\KM}_i=\bar{\KP}_i=:\bar K_i$ for the
common image of the torus generators.  We construct mutually inverse maps.

\emph{The map $p:\U_q(\sl(n))\to Q$.}  By
Theorem~\ref{thm:presentation}, $Q$ is generated by
$\bar K_a^{\pm1}$, $E_i$, $F_i$ with the relations obtained from
\eqref{eq:D1}--\eqref{eq:D5} by identifying $\KM_a=\KP_a$.  These are
precisely the relations \eqref{eq:Uq-torus}--\eqref{eq:SerreF} of
$\U_q(\sl(n))$: for \eqref{eq:Uq-torus}--\eqref{eq:SerreF} this is immediate;
the two conjugation relations of \eqref{eq:D2} and \eqref{eq:D4} become the
same relation $\bar K_aE_i\bar K_a^{-1}=q^{(\al_i)_a}E_i$ and
$\bar K_aF_i\bar K_a^{-1}=q^{-(\al_i)_a}F_i$; and \eqref{eq:D5} becomes
\[
E_iF_j-F_jE_i=\delta_{ij}\frac{\bar K_i\bar K_{i+1}^{-1}
-\bar K_i^{-1}\bar K_{i+1}}{q-q^{-1}}
=\delta_{ij}\frac{K_i-K_i^{-1}}{q-q^{-1}},
\]
with $K_i=\bar K_i\bar K_{i+1}^{-1}$, because in the quotient
$\kappa_i^{-,-1}=(\KM_i\KM_{i+1}^{-1})^{-1}
=\KP_i^{-1}\KP_{i+1}=K_i^{-1}$.  Hence the universal property of
$\U_q(\sl(n))$ gives a surjection $p:\U_q(\sl(n))\twoheadrightarrow Q$ with
$p(\hat K_a)=\bar K_a$, $p(E_i)=E_i$, $p(F_i)=F_i$.

\emph{The map $j:Q\to\U_q(\sl(n))$.}  Define
$j(\KM_a)=j(\KP_a)=\hat K_a$, $j(E_i)=E_i$, $j(F_i)=F_i$.  This respects all
relations \eqref{eq:D1}--\eqref{eq:D5}: for \eqref{eq:D1}--\eqref{eq:D3} this
is clear, for \eqref{eq:D4} because
$\hat K_aE_i\hat K_a^{-1}=q^{(\al_i)_a}E_i$ and
$\hat K_aF_i\hat K_a^{-1}=q^{-(\al_i)_a}F_i$ in $\U_q(\sl(n))$, and for
\eqref{eq:D5} because in $\U_q(\sl(n))$
\[
E_iF_j-F_jE_i=\delta_{ij}\frac{K_i-K_i^{-1}}{q-q^{-1}}
=\delta_{ij}\frac{\kappa_i^+-\kappa_i^{-,-1}}{q-q^{-1}},
\]
with the obvious meaning of $\kappa_i^\pm$ in $\U_q(\sl(n))$ (here
$\kappa_i^{+}=\kappa_i^{-}=\hat K_i\hat K_{i+1}^{-1}$, so
$\kappa_i^+-\kappa_i^{-,-1}=K_i-K_i^{-1}$).  Moreover $j$ annihilates the
ideal $(\KM_i-\KP_i)$, so it descends to $j:Q\to\U_q(\sl(n))$.

Finally $p$ and $j$ are mutually inverse on generators, hence mutually
inverse isomorphisms.
\end{proof}

\begin{proof}[Proof of Theorem~\ref{thm:B}]
By Theorem~\ref{thm:A} and Lemma~\ref{lem:quotient-ideal},
\[
\mathscr S_n(\mathfrak B^*)/I\ \cong\
D_n/(\KM_i\KP_i^{-1}-1,\ 1\le i\le n),
\]
and by Theorem~\ref{thm:quotient} the latter is isomorphic to
$\U_q(\sl(n))$.  The composite isomorphism sends the class of an arc to its
image in $\U_q(\sl(n))$ obtained from \eqref{eq:Phi}--\eqref{eq:Phi2}
followed by $j$; in particular the diagonal arcs map to the torus
$\hat K_i$, the simple lower arcs $x_i$ to the $F_i$, and the simple upper
arcs $y_i$ to the $E_i$ of the simply-connected quantum group.
\end{proof}

\begin{corollary}
The simply-connected quantum group $\U_q(\sl(n))=U_q^{\mathrm{ext}}(\sl(n))$
is generated by the images of the horizontal arcs, subject only to the
relations induced from $\mathscr S_n(\mathfrak B^*)$ and the gluing
relations $b_{ii}c_{ii}=1$.
\end{corollary}

%=====================================================================
\section{Concluding remarks}
%=====================================================================

\begin{remark}
Theorem~\ref{thm:quotient} is the simply-connected version of the classical
Drinfeld--Tanisaki realization of a quantized enveloping algebra as a
quotient of the double of its two Borels: the extra relations
$\KM_i=\KP_i$ identify the two copies of the simply-connected torus, and the
offset $1/n$ in the skew pairing guarantees (Lemma~\ref{thm:pairing-exists})
that the pairing descends to the extended, simply-connected Borels in the
first place.
\end{remark}

\begin{remark}
The map $\Phi$ of Theorem~\ref{thm:A} is ``explicit'' in the strongest sense:
each generator $b_{ij}$, $c_{ij}$ is sent to a single ordered monomial in the
Chevalley generators of $D_n$, and the inverse map $\Psi$ sends each
Chevalley generator to a single monomial in the arcs.  Consequently every
element of the double (and hence, via Theorem~\ref{thm:B}, of the quantum
group) acquires a unique normal form as a linear combination of ordered
monomials in the arcs, read off directly from the PBW basis
(Lemma~\ref{lem:PBW}) through \eqref{eq:Phi}--\eqref{eq:Phi2}.
\end{remark}

\begin{thebibliography}{10}

\bibitem{KS}
A.~U. Klimyk, K.~Schm\"udgen,
\emph{Quantum Groups and Their Representations},
Springer-Verlag, Berlin, 1997.

\bibitem{Tan}
T. Tanisaki,
\emph{Killing forms, Harish-Chandra isomorphisms, and universal
$R$-matrices for quantum algebras},
Invent.\ Math.\ \textbf{102} (1990), 541--566.

\bibitem{Jan}
J.~C. Jantzen,
\emph{Lectures on Quantum Groups},
Graduate Studies in Mathematics, vol.~6, American Mathematical Society,
Providence, RI, 1996.

\bibitem{FRT}
L.~D. Faddeev, N.~Yu. Reshetikhin, L.~A. Takhtajan,
\emph{Quantization of Lie groups and Lie algebras},
in: Algebraic Analysis, Vol.~I, Academic Press, Boston, 1988, pp.~129--139.

\bibitem{Maj}
S. Majid,
\emph{Foundations of Quantum Group Theory},
Cambridge University Press, Cambridge, 1995.

\bibitem{Lus}
G. Lusztig,
\emph{Introduction to Quantum Groups},
Progress in Mathematics, vol.~110, Birkh\"auser, Boston, 1993.

\end{thebibliography}

\end{document}
```

**Summary of the solution.** The document above proves the two requested statements in the following precise form.

* **Setup (\S1–\S2).** The simply-connected quantum group $\mathcal U_q(\mathfrak{sl}(n))=U_q^{\mathrm{ext}}(\mathfrak{sl}(n))$ of Klimyk–Schmüdgen, its simply-connected Borels $\mathcal U_q(\mathfrak b_\pm)$, and the Drinfeld–Tanisaki skew pairing with the prescribed values $\langle \hat K_i,\hat K_j\rangle=q^{1/n-\delta_{ij}}$, $\langle E_i,F_j\rangle=\delta_{ij}/(q^{-1}-q)$ are fixed; the $1/n$-offset is shown to be exactly what makes the pairing descend to the simply-connected (extended) Borels.

* **The double (\S3).** Following Tanisaki and Klimyk–Schmüdgen, the generalized quantum double $\mathcal U_q(\mathfrak b_-)\bowtie\mathcal U_q(\mathfrak b_+)$ is defined via the cross relations (T). The key structural result, Theorem 4.2, computes these cross relations on the generators and shows that the double has the explicit presentation with two commuting simply-connected tori, the standard Borel relations, the cross conjugation relations $\hat K_a^-E_i=q^{(\alpha_i)_a}E_i\hat K_a^-$, $\hat K_a^+F_i=q^{-(\alpha_i)_a}F_i\hat K_a^+$, and the mixed relation $E_iF_j-F_jE_i=\delta_{ij}\frac{\kappa_i^+-\kappa_i^{-,-1}}{q-q^{-1}}$; the Hopf structure is verified directly.

* **The isomorphism (\S4).** The bigon algebra $\mathcal S_n(\mathfrak B^*)$ of horizontal arcs is defined by generators $b_{ij}$ (lower arcs, $i\ge j$) and $c_{ij}$ (upper arcs, $i\le j$) and the arc relations (S1)–(S5). The explicit map $\Phi$ of Theorem 1.1,
$b_{ii}\mapsto\hat K_i^-$, $b_{ij}\mapsto\hat K_i^-\hat K_j^{-1}F_j\cdots F_{i-1}$, $c_{ii}\mapsto\hat K_i^{+,-1}$, $c_{ij}\mapsto\hat K_i^{+,-1}\hat K_j^+E_i\cdots E_{j-1}$,
together with its explicit inverse (diagonal arcs $\leftrightarrow$ tori, simple arcs $\leftrightarrow$ Chevalley generators) is shown to be a well-defined, mutually inverse pair of algebra homomorphisms, which proves problem (1).

* **The quotient (\S5).** Since $\Phi(b_{ii}c_{ii}-1)=\hat K_i^-\hat K_i^{+,-1}-1$, the quotient of the double by the diagonal torus identification is computed (Theorem 5.2) and shown to be exactly $\mathcal U_q(\mathfrak{sl}(n))$, because after identifying the two tori the cross relations collapse to the standard defining relations of the simply-connected quantum group. Combining the two results yields problem (2): $\mathcal U_q(\mathfrak{sl}(n))\cong\mathcal S_n(\mathfrak B^*)/(b_{ii}c_{ii}-1)$.